\chapter{绪论}
\label{chap:introduction}
\section{模板目的与使用范围}
本模板依据项目中提供的贵州大学研究生学位论文格式文件制作，覆盖硕士论文的前置页面、正文与结尾部分。
本示例中的文字、表格和图形仅用于演示格式，不代表真实研究结果。正式提交前须替换论文信息、摘要与成果示例，并完成声明签名。

正文使用五号宋体，段首缩进两个汉字，段落之间不额外增加空白。章标题使用小四号黑体，节标题使用五号黑体；章、节编号左对齐，编号与标题之间空一个字。每章另起一页，双面模式下正文的第一页从右页开始。

\section{文献引用与写作约定}
顺序编码制示例使用学校附件列出的图书\cite{zhang2002}和期刊论文\cite{li2000}。
参考文献由数据库自动生成，顺序与正文第一次引用的顺序一致。这里仅保留两条真实书目信息演示流程，正式论文须根据实际引用补足文献；学校规定硕士不少于 35 篇。

\subsection{交叉引用}
正文应先介绍图、表或公式，再放入对应内容。使用标签与交叉引用可避免章节调整后手动编号出错。
第~\ref{chap:layout}~章展示图表、数学公式和跨页表格，附录~\ref{app:extra}~给出补充材料示例。

\subsection{注释}
需要解释、但没有具体文献来源的内容可使用脚注。脚注位于当前页底部\footnote{这是注释示例。学校建议控制注释数量；本模板默认采用普通阿拉伯数字脚注编号。}。
引用他人成果时应使用文献引用命令，而不是用脚注代替文献条目。

\section{默认值与院系要求}
附件没有明确给出页边距、正文行距或页眉样式。本模板默认使用上、下和外侧 25 mm、内侧 30 mm 的页边距，正文基线间距为 20 bp；页脚居中显示页码，不设置页眉。这些是可修改的排版默认值，具体使用时以院系补充要求为准。

默认字体采用 TeX Live 自带的 Fandol 宋体、黑体和楷体，便于在 Linux、Windows 和在线环境编译。若院系要求中易宋体，应安装合法字体并切换字体配置，不能将字体替代方案当成指定字体。
